% !TEX program = xelatex
% !BIB program = biber
%-------------------------------------------------------------------------------
% CONFIGURATIONS
%-------------------------------------------------------------------------------
% A4 paper size by default, use 'letterpaper' for US letter
\documentclass[11pt, a4paper]{russell}
\usepackage{hyperref}

% Iterate publications.bib and filter by entry type (no per-entry listing).
\usepackage[backend=biber, sorting=none, style=numeric, defernumbers=true]{biblatex}
\addbibresource{publications.bib}

% Drop the default quotes around titles in the numeric style.
\DeclareFieldFormat{title}{#1}
\DeclareFieldFormat[article,inbook,incollection,inproceedings,patent,thesis,unpublished]{title}{#1}

% Hanging-indent entries with no [N] label — matches the prior \publication look.
% Trick: leftmargin pushes the whole entry right, itemindent pulls line 1 back.
\defbibenvironment{bibliography}
  {\list{}
     {\setlength{\leftmargin}{1.5em}%
      \setlength{\itemindent}{-1.5em}%
      \setlength{\labelwidth}{0pt}%
      \setlength{\labelsep}{0pt}%
      \setlength{\itemsep}{\smallskipamount}%
      \setlength{\parsep}{0pt}}}
  {\endlist}
  {\item}
% Configure page margins with geometry
\geometry{left=1.4cm, top=.8cm, right=1.4cm, bottom=1.5cm, footskip=.5cm}

% Fix overlapping text by increasing section/entry spacing
\renewcommand{\acvSectionTopSkip}{3.5mm}
\renewcommand{\acvSectionContentTopSkip}{2.5mm}

% Specify the location of the included fonts
\fontdir[fonts/]

% Color for highlights
% russell Colors: russell-emerald, russell-skyblue, russell-red, russell-pink, russell-orange
%                 russell-nephritis, russell-concrete, russell-darknight, russell-purple
\colorlet{russell}{russell-black}
% Uncomment if you would like to specify your own color
% \definecolor{russell}{HTML}{CA63A8}

% Colors for text
% Uncomment if you would like to specify your own color
% \definecolor{darktext}{HTML}{414141}
% \definecolor{text}{HTML}{333333}
% \definecolor{graytext}{HTML}{5D5D5D}
% \definecolor{lighttext}{HTML}{999999}

% Set false if you don't want to highlight section with russell color
\setbool{acvSectionColorHighlight}{true}

% If you would like to change the social information separator from a pipe (|) to something else
\renewcommand{\acvHeaderSocialSep}{\quad\textbar\quad}
\newcommand{\githubURL}[1]{\href{https://#1}{\faGithubSquare\acvHeaderIconSep#1}}
\newcommand{\paperURL}[2]{\href{https://#1}{\faFile*[regular]\acvHeaderIconSep#2}}

%-------------------------------------------------------------------------------
%	PERSONAL INFORMATION
%	Comment any of the lines below if they are not required
%-------------------------------------------------------------------------------
% Available options: circle|rectangle,edge/noedge,left/right
% \photo[rectangle,edge,right]{./examples/profile}
\name{El Hajj Chehade,}{Saiid}
% \position{Software Architect{\enskip\cdotp\enskip}Security Expert}
\address{BC 262, EPFL CH-1014 Lausanne, Switzerland}

\email{saiid.elhajjchehade@epfl.ch}
%\dateofbirth{January 1st, 1970}
\homepage{saiid.ch}
\github{github.com/saiid2001}
\linkedin{linkedin.com/in/saiid-hc}
% \gitlab{gitlab-id}
% \stackoverflow{SO-id}{SO-name}
% \twitter{@twit}
% \skype{skype-id}
% \reddit{reddit-id}
% \medium{madium-id}
% \kaggle{kaggle-id}
% \googlescholar{googlescholar-id}{name-to-display}
%% \firstname and \lastname will be used
\googlescholar{gF0rvJAAAAAJ}{}
% \extrainfo{extra information}

% \quote{``Simplicity is deceptively complicated."}


%-------------------------------------------------------------------------------
\begin{document}

\nocite{*}
% Print the header with above personal informations
% Give optional argument to change alignment(C: center, L: left, R: right)
\makecvheader
% Print the footer with 3 arguments(<left>, <center>, <right>)
% Leave any of these blank if they are not needed
\makecvfooter
  {\today}
  {}
  {\thepage}


%-------------------------------------------------------------------------------
%	CV/RESUME CONTENT
%	Each section is imported separately, open each file in turn to modify content
%-------------------------------------------------------------------------------

%-------------------------------------------------------------------------------
%	SECTION TITLE
%-------------------------------------------------------------------------------
\cvsection{Education}


%-------------------------------------------------------------------------------
%	CONTENT
%-------------------------------------------------------------------------------
\begin{cventries}

%---------------------------------------------------------
  \cventry
    {PhD in Computer Science} % Degree
    {Ecole Polytechnique Fédérale de Lausanne (EPFL)} % Institution
    {Lausanne, Switzerland} % Location
    {Sep 2023 -- Present} % Date(s)
    {
      \begin{cvitems} % Description(s) bullet points
        \item {Security and Privacy Field. Specializing in Web and Browser Privacy and Security.}
        \item {EDIC Fellowship Recipient.}
        \item {\textbf{Advisor:} Prof. Carmela Troncoso.}
      \end{cvitems}
    }

  \cventry
    {B.E. in Computer and Communication Engineering} % Degree
    {American University of Beirut} % Institution
    {Beirut, Lebanon} % Location
    {Aug 2019 -- Jun 2023} % Date(s)
    {
      \begin{cvitems} % Description(s) bullet points
        \item {Minor in Applied Mathematics. Graduated with High Distinction (GPA: 4.23/4.0).}
        \item {Dean's Honor List (all semesters), Merit Scholarship (3rd in Lebanese national secondary exams).}
        \item {\textbf{Final Year Project:} Decentralized private polling system using threshold homomorphic encryption and touch-biometric authentication.}
      \end{cvitems}
    }
%---------------------------------------------------------
\end{cventries}

%-------------------------------------------------------------------------------
%	SECTION TITLE
%-------------------------------------------------------------------------------


\cvsection{Publications}

%-------------------------------------------------------------------------------
%	CONTENT
%-------------------------------------------------------------------------------

\publication{vsf2025}
    \githubURL{github.com/Saiid2001/vsf}
    
\publication{flfp2025}
    \githubURL{github.com/spring-epfl/flfp}

\publication{sinbad2024}
    \githubURL{github.com/spring-epfl/sinbad}

\subsubsection*{Work in Progress / Submitted}

\cvsection{Talks \& Presentations}

\begin{cvhonors}

   
  \cvhonor
    {Workshop Talk} % Type
    {(Planned) Breaking the web is good for privacy, HotPET } % Title, Venue
    {Calgary, Canada} % Location
    {2026} % Date

  \cvhonor
    {Workshop Talk} % Type
    {(Planned) Semantic Web Page Representations for Large-Scale LLM-Powered Web Measurement, Dagstuhl Seminar On Web Application Security  } % Title, Venue
    {Dagstuhl, Germany} % Location
    {2026} % Date
 
  \cvhonor
    {Workshop Talk} % Type
    {\href{https://www.youtube.com/watch?v=vEHYKxPq7yE}{Stealthy Fingerprinting Users with personalized blocking rules}, Ad-Filtering Dev Summit} % Title, Venue
    {Limassol, Cyprus} % Location
    {2025} % Date

  \cvhonor
    {Invited Talk}
    {Customized Ad-blocker fingerprinting,  New York University (NYU)}
    {Brooklyn, NY, USA}
    {2025}

  \cvhonor
    {Invited Talk}
    {Browser Vendor Fingerprinting, Max-Plank Institute for Security and Privacy (MPI-SP)}
    {Bochum, Germany}
    {2025}

  \cvhonor
    {Invited Talk}
    {Web Breakage Detection, DuckDuckGo Research Team}
    {Remote}
    {2024}

  \cvhonor
    {Invited Talk}
    {Web Breakage Detection, Brave Browser Research Team}
    {Remote}
    {2024}

\end{cvhonors}

%-------------------------------------------------------------------------------
%	SECTION TITLE
%-------------------------------------------------------------------------------
\cvsection{Experience}


%-------------------------------------------------------------------------------
%	CONTENT
%-------------------------------------------------------------------------------
\begin{cventries}

%---------------------------------------------------------

\cventry
    {Brave Browser and Imperial College London} % Organisation
    {Research Intern} % Job title
    {London, United Kingdom} % Location
    {June.  2025 - August. 2025} % Date(s)
    {
      \begin{cvitems} % Description(s) of tasks/responsibilities
        \item {Worked with Prof. Hamed Haddadi and the Brave research team on creating optimized page representations for text-based web agents to be used in large-scale security and privacy web measurements.}
        \item {Produced a middle layer between DOM and Accessibility tree that preserves the semantic meaning in the page.}
        \item {Generated proposal and validated findings during internship, and continuing with an upcoming academic submission.}
      \end{cvitems}
    }
    
\cventry
    {Secure Web Application Group - CISPA} % Organisation
    {Visiting Researcher} % Job title
    {St. Ingbert, Germany} % Location
    {Jul.  2024 - Sept. 2024} % Date(s)
    {
      \begin{cvitems} % Description(s) of tasks/responsibilities
        \item {Worked with Prof. Ben Stock on unsafe parameters and authorization issues in HTTP requests for authenticated sessions.}
        \item {Developed a mirrored browser setup for synchronized requests from two controlled user accounts.}
        \item {Managed a group of 5 volunteers for large-scale semi-manual experiments.}
      \end{cvitems}
    }


\cventry
    {SPRING Lab - EPFL} % Organisation
    {Research Assistant} % Job title
    {Lausanne, Switzerland} % Location
    {Sept.  2023 - Present.} % Date(s)
    {
      \begin{cvitems} % Description(s) of tasks/responsibilities
        \item {Working on the intersection between web-measurement, web-privacy and browsers.}
      \end{cvitems}
    }

 \cventry
    {Test Prep Institute} % Organisation
    {Lead Developer (Full Stack)} % Job title
    {Beirut, Lebanon} % Location
    {May.  2022 - Sept. 2023} % Date(s)
    {
      \begin{cvitems} % Description(s) of tasks/responsibilities
        \item {Managed team of 3 to build an online educational platform. Helped company implement automated tests and decrease student processing time by 70\%.}
        \item {Helped design an architecture spanning a secure student assessment environment, adaptive questions based on performance, classroom database system, user access control and live tracking exam statistics. These features allowed the CEO to expand to a larger market of institutions in Lebanon.}
        \item {Managed \textit{Git} repository and Deployment on \textit{Heroku} for \textit{Django} Build with embedded \textit{React} frontend through \i{Webpack}. This helped get 98\% up-time on the server.}
        \item {Lead strategy development by setting task priorities, and timelines. This helped deploy the initial version in one month.}
        \item {Got offer to stay for equity to create the IT department.}
      \end{cvitems}
    }
    
  \cventry
    {American University of Beirut} % Organisation
    {Undergraduate Research Assistant} % Job title
    {Beirut, Lebanon} % Location
    {Sept.  2022 - Jan. 2023} % Date(s)
    {
      \begin{cvitems} % Description(s) of tasks/responsibilities
        \item {Worked alongside Prof. Ibrahim Issa on the problem of "Fairness in Machine Learning".}
        \item {Formulated the problem statement and the direction of the project.}
        \item {Compiled literature review to have an understanding of the research domain standing.}
        \item {Prepared a report detailing the taxonomy, directions, and a new definition for Fairness.}
        \item { \textbf {Academic areas:} Differential privacy, statistics \& machine learning, linear algebra, optimization.}
      \end{cvitems}
    }

%---------------------------------------------------------
  \cventry
    {SPRING Lab EPFL} % Organisation
    {Research Intern} % Job title
    {Lausanne, Switzerland} % Location
    {Jun.  2022 - Aug. 2022} % Date(s)
    {
      \begin{cvitems} % Description(s) of tasks/responsibilities
        \item Presented a differential DOM tree comparison algorithm and redefined web-breakage problem definition that paved the way for a promising paper endorsed by the lab advisor Prof. Carmela Troncoso.
        \item Built a pipeline to classify web-page breakages to increase the usability of Ad-blockers: Feature selection, and random-forests. 
        \item Designed crawl experiments and evaluated model. I achieved comparable accuracy to state-of-the-art techniques.
        \item Published a USENIX Artifact for WebGraph: Smart Ad classifier under supervision of Dr. Sandra Siby.
        \item Published the first online Ad-block enabled web crawling tool that was recognized by my supervisor.\\\textit{ GitHub: \href{https://www.github.com/Saiid2001/OpenWPM}{Saiid2001/OpenWPM}.} 
        \item Created GitHub Crawlers to automate the ab-lock history tracking in Python.
        \item { \textbf{Technical skills: } python, OpenWPM, sklearn, multiprocessing, sqlite, linux, Git.}
        \item { \textbf {Academic areas:} Web-crawling, machine learning, networking.}
      \end{cvitems}
    }

%---------------------------------------------------------

  \cventry
    {NiftyCraft} % Organisation
    {Block-chain Developer Part-time} % Job title
    {Remote, UAE} % Location
    {Jan. 2022 - Aug. 2022} % Date(s)
    {
      \begin{cvitems} % Description(s) of tasks/responsibilities
        \item Published a smart-contract for NFT tokens for in-game usage.
        \item Researched and Integrated NFT authentication for exclusive content serving for token holders. Deployed and running in production. 
        \item { \textbf{Technical skills: } Solidity, reactJs,  Git.}
        \item { \textbf {Academic areas:} cryptography}
      \end{cvitems}
    }

%---------------------------------------------------------


  \cventry
    {American University of Beirut} % Organisation
    {Research Assistant} % Job title
    {Beirut, Lebanon} % Location
    {Sep.  2020 - Dec. 2021} % Date(s)
    {
      \begin{cvitems} % Description(s) of tasks/responsibilities
        \item Led a Human Centered Design (HCD) research project with Prof. Hazem Hajj to solve students’ social needs in online education.
        \item \href{https://drive.google.com/file/d/1yCXrCYoeaBECwf5p4byoiabFzedO1Zgr/view?usp=sharing}{Developed a full virtual campus platform on ElectronJS/Flask in the span of 1 year}.
        \item Created a mesh topography conference interface using WEBrtc for up to 6 users.
        \item Scraped University website for courses information using Python.
        % \item Rewired Bitmoji API to create custom avatars.
        \item { \textbf{Technical skills: } WebRTC, ReactJs, javascript, flask, python, docker, Git.}
        \item { \textbf {Academic areas:} human-centered design, communication, psychology}
      \end{cvitems}
    }

%---------------------------------------------------------


  % \cventry
  %   {Sila.Ai Startup} % Organisation
  %   {Full-stack Developer} % Job title
  %   {Beirut, Lebanon} % Location
  %   {Sep. 2020 - Mar. 2021} % Date(s)
  %   {
  %     % \begin{cvitems} % Description(s) of tasks/responsibilities
  %     %   \item Coded Python algorithms to parse time series data on online exam platform.
  %     %   \item Visualized student activity and predicted malicious activity from time-data.
  %     %   \item Developed online exam platform (Frontend JavaScript, HTML, CSS) and backend with Docker, Flask and MongoDB.
  %     % \end{cvitems}
  %   }


%---------------------------------------------------------

% \cventry
%     {Falafel Games} % Organisation
%     {Game Development Intern} % Job title
%     {Remote, UAE} % Location
%     {Aug. 2020 - Sep. 2020} % Date(s)
%     { }
%---------------------------------------------------------

\end{cventries}

\cvsection{Teaching}

%-------------------------------------------------------------------------------
%	CONTENT
%-------------------------------------------------------------------------------
\begin{cvhonors}

  %---------------------------------------------------------
  \cvhonor
  {Fall '23, '24, '25} % Award
  {Teaching Assistant for Introduction to Computer Security} % Event
  {EPFL, Switzerland} % Location
  {COM-301} % Date(s)
  \cvhonor
  {Spring '25} % Award
  {Teaching Assistant for Advanced Topics in Privacy Enhancing Technologies} % Event
  {EPFL, Switzerland} % Location
  {CS-523} % Date(s)

\end{cvhonors}

\cvsection{Supervised Students}

%-------------------------------------------------------------------------------
%	CONTENT
%-------------------------------------------------------------------------------
\begin{cvhonors}

  %---------------------------------------------------------

  \cvhonor
  {Master Semester Project}
  {Cristiano Sartori, Synthetic Browser Personas and PET Effects on Ad Distributions}
  {EPFL, Switzerland}
  {Spring '26}

  \cvhonor
  {Master Semester Project}
  {Alexandre Dunoyer de Segonzac, Deep Advertisement Scanning and Clustering}
  {EPFL, Switzerland}
  {Fall '25}

  \cvhonor
  {Master Semester Project} % Award
  {Julien Mettler, Systematic evaluation of logged-in cookie sensitivity} % Event
  {EPFL, Switzerland} % Location
  {Spring '24} % Date(s)

  \cvhonor
  {Master Semester Project} % Award
  {Jonas Gloning, Evaluating the privacy impact of Chrome V3 on web extensions} % Event
  {EPFL, Switzerland} % Location
  {Spring '24} % Date(s)

  \cvhonor
  {Master Semester Project} % Award
  {Pierre-Louis Bravais, Automated agents for web measurements} % Event
  {EPFL, Switzerland} % Location
  {Fall '24} % Date(s)

  \cvhonor
  {Master Semester Project} % Award
  {Siim Markus Marvet, Investigating Server-side Google Tag Manager} % Event
  {EPFL, Switzerland} % Location
  {Fall '24} % Date(s)

\end{cvhonors}

%-------------------------------------------------------------------------------
%	SECTION TITLE
%-------------------------------------------------------------------------------
\newpage

\cvsection{Awards \& Honors}



%-------------------------------------------------------------------------------
%	CONTENT
%-------------------------------------------------------------------------------
\begin{cvhonors}

%---------------------------------------------------------
  \cvhonor
    {EDIC Fellowship} % Award
    {School of Computer Science, EPFL} % Event
    {Switzerland} % Location
    {2023} % Date(s)
  \cvhonor
    {Fellowship} % Award
    {Summer Research Institute} % Event
    {Switzerland} % Location
    {2023} % Date(s)
  \cvhonor
    {7x Dean's Honor List} % Award
    {American University of Beirut} % Event
    {Lebanon} % Location
    {2019--2023} % Date(s)
  \cvhonor
    {AUB / CNRS Merit Scholarship} % Award
    {3rd rank in Lebanese national secondary education exams} % Event
    {Lebanon} % Location
    {2019} % Date(s)
  \cvhonor
    {Prime Minister Honor Award} % Award
    {Awarded for academic excellence in secondary education} % Event
    {Lebanon} % Location
    {2019} % Date(s)
\end{cvhonors}

%-------------------------------------------------------------------------------
%	SECTION TITLE
%-------------------------------------------------------------------------------
\cvsection{Outreach}

%-------------------------------------------------------------------------------
%	CONTENT
%-------------------------------------------------------------------------------
\begin{cventries}

%---------------------------------------------------------
  \cventry
    {EPFL} % Organisation
    {RAMP Program} % Project
    {Lausanne, Switzerland} % Location
    {Present} % Date(s)
    {
      Volunteer as Mentor for PhD applicants at EPFL. Responsible for reviewing and giving feedback on CVs and Statement of Purpose letters.
    }


%---------------------------------------------------------
  \cventry
    {Robotics Club, American University of Beirut} % Organisation
    {Linux Workshop} % Project
    {Beirut, Lebanon} % Location
    {Nov. 2022} % Date(s)
    {
      \begin{cvitems} % Description(s) of project
        \item Volunteered with the Robotics Club at AUB to deliver an introductory Linux workshop.
        \item Presented in front of 50 students. They demanded a second one to be given soon.
      \end{cvitems}
    }

%---------------------------------------------------------
  \cventry
    {Independent Initiative} % Organisation
    {Kids Robotics Workshop} % Project
    {Chehim, Lebanon} % Location
    { } % Date(s)
    {
      \begin{cvitems} % Description(s) of project
        \item Organized local robotics training course for 20 children. I motivated student interest for robotics and computer related fields of study.
         \item Prepared curriculum, book, and presentations for the course.
      \end{cvitems}
    }


%---------------------------------------------------------
  \cventry
    {"Popular Rescue" local organization } % Organisation
    {First-Aid training} % Project
    {Chehim, Lebanon} % Location
    {Jul. 2018} % Date(s)
    { }

%---------------------------------------------------------
\end{cventries}

%-------------------------------------------------------------------------------
%	SECTION TITLE
%-------------------------------------------------------------------------------

\cvsection{Skills \& Languages}
%-------------------------------------------------------------------------------
%	CONTENT
%-------------------------------------------------------------------------------
\begin{cvskills}

  %---------------------------------------------------------
  \cvskill
  {Programming} % Category
  {Python (Pandas, TensorFlow, NumPy, Scikit-learn, Spacy), C/C++, JavaScript, SQL, Java, CUDA, Django} % Skills

  %---------------------------------------------------------
  \cvskill
  {Tools} % Category
  {Linux, Bash, \LaTeX, React, Docker, Git, Figma, LXD, LLM APIs, Playwright} % Skills

  %---------------------------------------------------------
  \cvskill
  {Languages} % Category
  {English (Professional, IELTS 8.5), French (Bilingual, DELF B2), Arabic (Native)} % Skills

\end{cvskills}


%-------------------------------------------------------------------------------
\end{document}
